% !TeX root = ../概率论与数理统计（缪柏其）笔记.tex
% 由主文档通过 \input 载入；请编译主文档。
\setcounter{chapter}{5} % 保留原书章号。
\chapter{参数点估计}
\label{chap:study-point-estimation}

从已知总体分布推导样本的性质，是概率论的基本方向；从样本反推分布中的未知参数，则进入统计推断。本章以教材第六章为线索，讨论如何构造点估计、如何比较估计量，以及样本量增大后估计误差怎样变化。构造方法与评价准则是两个不同问题：矩估计或最大似然估计说明估计量怎样产生，无偏性、均方误差和相合性等说明它在重复抽样中有什么性质。
% 教材依据：缪柏其、张伟平《概率论与数理统计》（2022），第六章，正文第 207--229 页，习题第 230--239 页。

\section{参数点估计的概念}
\label{sec:study-point-estimation-concept}

设总体分布属于分布族 $\{F_\theta:\theta\in\Theta\}$，其中 $\theta=(\theta_1,\ldots,\theta_k)$ 是未知参数，$\Theta\subseteq\mathbb R^k$ 称为\emph{参数空间（parameter space）}。本章除特别说明外，$X_1,\ldots,X_n$ 均为来自该总体的简单随机样本，即独立且同分布；$\boldsymbol X=(X_1,\ldots,X_n)$ 表示随机样本，$\boldsymbol x=(x_1,\ldots,x_n)$ 表示它的一次观测值。记 $\mathbb E_\theta$、$\operatorname{Var}_\theta$ 和 $\mathbb P_\theta$ 为参数固定在 $\theta$ 时的期望、方差和概率。\emph[caution]{频率学派的点估计把参数视为固定的未知常数，随机性来自抽样。}

\begin{definition}[点估计量与估计值]
\label{def:study-point-estimator}
设 $g:\Theta\to\mathbb R$ 为已知函数。用统计量
\[
  T_n=T_n(\boldsymbol X)
\]
估计 $g(\theta)$，称 $T_n$ 为 $g(\theta)$ 的\emph{点估计量（point estimator）}；观测到 $\boldsymbol x$ 后所得的数 $T_n(\boldsymbol x)$ 称为相应的\emph{估计值（estimate）}。统计量的计算规则不能含有未知参数。取 $g(\theta)=\theta_j$ 就是在估计某个参数分量，向量参数也可逐分量估计。
\end{definition}

例如，总体具有有限均值 $\mu$ 时，样本均值 $\overline X=n^{-1}\sum_{i=1}^nX_i$ 是 $\mu$ 的一个估计量；测得三个寿命值 $5,6.4,5.9$ 后，估计值为 $\overline x=17.3/3\approx5.77$。前者具有抽样分布，后者是当前数据算出的数。点估计提供一个代表值，并不同时给出误差范围；误差的量化要结合估计量的分布，进一步引出下一章的区间估计。

\begin{example}
\label{ex:study-voltage-estimators}
教材例 6.1 的 $20$ 个击穿电压观测值之和为 $555.86$，最小值、最大值分别为 $24.46,30.88$，中间两个次序值为 $27.94,27.98$，单位均为 $\mathrm{kV}$。若用正态总体建模，怎样估计总体均值 $\mu$？
\end{example}

\begin{solve}
正态分布的均值、中位数和对称中心相同，因此可考虑样本均值、样本中位数（sample median）和中程数（midrange）。令 $X_{(1)}\le\cdots\le X_{(n)}$ 为由小到大排列的次序统计量，三种估计值分别为
\[
  \overline x=\frac{555.86}{20}=27.793,\qquad
  \widetilde x=\frac{27.94+27.98}{2}=27.960,\qquad
  \frac{x_{(1)}+x_{(20)}}2=27.670.
\]
第三种使用最大值与最小值的平均，不能与极差的一半 $(x_{(20)}-x_{(1)})/2$ 混淆。不同估计值不完全相同并不构成矛盾；比较估计方法，需要研究它们在重复抽样中的整体误差，而不能仅凭这一次数据判断谁离未知的 $\mu$ 更近。
\end{solve}

\section{矩估计法}
\label{sec:study-method-of-moments}

\subsection{以样本矩代替总体矩}

对整数 $r\ge1$，若 $\mathbb E_\theta|X|^r<\infty$，总体的 $r$ 阶原点矩和中心矩分别记为
\[
  \alpha_r(\theta)=\mathbb E_\theta(X^r),\qquad
  \mu_r(\theta)=\mathbb E_\theta[(X-\alpha_1(\theta))^r],
\]
相应的样本矩记为
\begin{equation}
  A_r=\frac1n\sum_{i=1}^nX_i^r,\qquad
  M_r=\frac1n\sum_{i=1}^n(X_i-\overline X)^r.
  \label{eq:study-sample-moments}
\end{equation}
原点矩可用大数定律得到 $A_r\xrightarrow{\mathbb P}\alpha_r(\theta)$；中心矩展开为样本原点矩的多项式后，也可得到 $M_r\xrightarrow{\mathbb P}\mu_r(\theta)$。这说明在相应矩存在时，样本矩能够近似总体矩。

\begin{definition}[矩估计]
\label{def:study-moment-estimator}
设选定的 $k$ 个总体矩为 $h_1(\theta),\ldots,h_k(\theta)$，对应样本矩为 $\widehat h_1,\ldots,\widehat h_k$。将矩之间的近似关系写成方程组
\begin{equation}
  h_j(\widehat\theta_{\mathrm M})=\widehat h_j,
  \qquad j=1,\ldots,k,
  \label{eq:study-moment-equations}
\end{equation}
若能在参数空间内选出一个由样本确定的解，就称它为参数的\emph{矩估计量（method-of-moments estimator）}。相应的参数函数估计为 $g(\widehat\theta_{\mathrm M})$。
\end{definition}

方程个数与参数个数相同并不保证解存在或唯一，求解后还须核对参数约束；矩存在也不保证所选矩能识别参数，例如对称分布的二阶矩可能无法区分某个位置参数的正负。矩估计无需写出完整密度，但必须知道选定总体矩与未知参数的关系。通常先使用低阶矩，因为高阶矩对尾部观测更敏感；这是一条方法选择的经验，不能脱离分布与误差准则断言低阶矩估计总是更优。

\subsection{常见总体与参数函数}

\begin{example}
\label{ex:study-normal-moments}
设 $X_i\sim N(\mu,\sigma^2)$，样本量 $n\ge2$，$\mu\in\mathbb R,\sigma^2>0$ 均未知，求二者的矩估计。
\end{example}

\begin{solve}
由 $\mathbb E(X)=\mu$、$\mathbb E(X^2)=\mu^2+\sigma^2$，矩方程给出
\begin{equation}
  \widehat\mu_{\mathrm M}=\overline X,\qquad
  \widehat{\sigma^2}_{\mathrm M}=A_2-\overline X^2
    =M_2=\frac1n\sum_{i=1}^n(X_i-\overline X)^2.
  \label{eq:study-normal-moment-estimates}
\end{equation}
当 $n\ge2$ 时，通常另记
\begin{equation}
  S^2=\frac1{n-1}\sum_{i=1}^n(X_i-\overline X)^2
     =\frac n{n-1}M_2,\qquad S=\sqrt{S^2}.
  \label{eq:study-unbiased-sample-variance}
\end{equation}
$S^2$ 是对 $M_2$ 作无偏修正后的估计，二者不能在推导中不加说明地互换。标准差 $\sigma$ 的直接矩估计为 $\sqrt{M_2}$，使用 $S$ 则是另一种选择。
\end{solve}

\begin{example}
\label{ex:study-rate-moments}
分别求指数分布、泊松分布和二项分布参数的低阶矩估计。
\end{example}

\begin{solve}
本章以速率参数定义指数分布：
\begin{equation}
  X\sim\operatorname{Exp}(\lambda)
  \quad\Longleftrightarrow\quad
  f(x;\lambda)=\lambda e^{-\lambda x}\mathbf1_{\{x\ge0\}},
  \qquad \lambda>0.
  \label{eq:study-exponential-rate}
\end{equation}
它的均值为 $1/\lambda$，所以 $\widehat\lambda_{\mathrm M}=1/\overline X$，而平均寿命 $1/\lambda$ 的矩估计为 $\overline X$。若 $X\sim\operatorname{Pois}(\lambda)$，由 $\mathbb E(X)=\lambda$ 得 $\widehat\lambda_{\mathrm M}=\overline X$。若 $X\sim\operatorname{Bin}(N,p)$，其中试验次数 $N\ge1$ 已知，则 $\mathbb E(X)=Np$，故 $\widehat p_{\mathrm M}=\overline X/N$；只有 $N=1$ 时才直接等于 $\overline X$。
\end{solve}

\begin{example}
\label{ex:study-uniform-moments}
设 $X_i\sim U(a,b)$，样本量 $n\ge2$，其中 $a<b$ 均未知，求矩估计。
\end{example}

\begin{solve}
由 $\mathbb E(X)=(a+b)/2$、$\operatorname{Var}(X)=(b-a)^2/12$，解方程
\[
  \frac{a+b}{2}=\overline X,\qquad
  \frac{(b-a)^2}{12}=M_2
\]
得到
\begin{equation}
  \widehat a_{\mathrm M}=\overline X-\sqrt{3M_2},\qquad
  \widehat b_{\mathrm M}=\overline X+\sqrt{3M_2}.
  \label{eq:study-uniform-moment-estimates}
\end{equation}
取正平方根是因为 $b-a>0$。当 $n\ge2$ 时，连续抽样下 $M_2>0$ 几乎必然成立。若按教材中的修正做法用 $S^2$ 替代 $M_2$，会得到 $\overline X\mp\sqrt3S$，这是相近但不同的估计。矩方程只匹配矩，并不保证估计出的区间包含全部观测值。
\end{solve}

参数函数可在得到参数估计后代入。例如，对于正态总体，设 $\Phi$ 为标准正态分布函数，则固定阈值 $c$ 对应的尾概率
\[
  q=\mathbb P(X>c)=1-\Phi\!\left(\frac{c-\mu}{\sigma}\right)
\]
的直接矩估计为 $1-\Phi((c-\overline X)/\sqrt{M_2})$，其中要求 $M_2>0$。若总体方差非零且四阶绝对矩有限，在 $M_2>0$ 时，偏度（skewness）与峰度（kurtosis）的矩估计分别为
\[
  \widehat\gamma_1=\frac{M_3}{M_2^{3/2}},\qquad
  \widehat\beta_2=\frac{M_4}{M_2^2};
\]
这里峰度采用未减去 $3$ 的约定，超额峰度的估计为 $\widehat\beta_2-3$。代入构造一般不保持无偏性。

\section{最大似然估计}
\label{sec:study-maximum-likelihood}

\subsection{固定样本，比较参数}

矩估计使用分布的若干数字特征，最大似然法则使用样本的联合分布。若总体概率密度或概率质量函数为 $f(x;\theta)$，独立样本在给定参数下的联合密度或联合概率质量为 $\prod_{i=1}^nf(x_i;\theta)$。

\begin{definition}[似然函数与最大似然估计]
\label{def:study-likelihood-mle}
固定观测值 $\boldsymbol x$，把联合密度或联合概率质量视为参数 $\theta$ 的函数，
\begin{equation}
  L(\theta;\boldsymbol x)=\prod_{i=1}^nf(x_i;\theta),
  \qquad\theta\in\Theta,
  \label{eq:study-likelihood}
\end{equation}
称为\emph{似然函数（likelihood function）}。若存在最大值点
\begin{equation}
  \widehat\theta_{\mathrm{ML}}(\boldsymbol x)
    \in\operatorname*{arg\,max}_{\theta\in\Theta}
       L(\theta;\boldsymbol x),
  \label{eq:study-mle-definition}
\end{equation}
则称其为\emph{最大似然估计（maximum likelihood estimate，MLE）}；把样本值换成随机样本便得到最大似然估计量。
\end{definition}

离散情形中，似然是不同参数下出现当前样本的概率；连续情形中，它是密度值，不能解释为取到某个精确样本点的正概率。\emph[caution]{似然函数不是参数的概率分布，也不要求对参数积分为 $1$。}它只在给定模型内比较参数对当前数据的支持程度，并不保证最大值对应真实参数。

当 $L>0$ 时，因对数严格递增，最大化 $L$ 等价于最大化\emph{对数似然（log-likelihood）}
\[
  \ell(\theta;\boldsymbol x)=\log L(\theta;\boldsymbol x)
      =\sum_{i=1}^n\log f(x_i;\theta).
\]
可导且位于参数空间内部的最大值点必须满足似然方程 $\partial\ell/\partial\theta_j=0$，但方程的根只是候选点，仍须检查整体最大性、边界和参数约束。若上确界只在参数空间外的边界趋近而未达到，则通常意义的 MLE 不存在；若存在多个最大值点，还须给出不含未知参数的选取规则。数值优化得到的驻点也不能自动当作全局最大值。

\begin{proposition}[最大似然估计的不变性]
\label{prop:study-mle-invariance}
若 $\widehat\theta_{\mathrm{ML}}$ 是 $\theta$ 的最大似然估计，则 $g(\widehat\theta_{\mathrm{ML}})$ 是参数函数 $g(\theta)$ 的最大似然估计。这里对 $\eta=g(\theta)$，以
\[
  L_g(\eta;\boldsymbol x)
    =\sup_{\theta:\,g(\theta)=\eta}L(\theta;\boldsymbol x),
  \qquad\eta\in g(\Theta),
\]
作为消去其余参数后的剖面似然（profile likelihood）。
\end{proposition}

\begin{proof}
对任意 $\eta$，都有 $L_g(\eta;\boldsymbol x)\le L(\widehat\theta_{\mathrm{ML}};\boldsymbol x)$；在 $\eta=g(\widehat\theta_{\mathrm{ML}})$ 处，取 $\theta=\widehat\theta_{\mathrm{ML}}$ 可达到右端。因此该 $\eta$ 为最大值点，不要求 $g$ 为一一映射。
\end{proof}

\subsection{光滑似然与参数边界}

\begin{example}
\label{ex:study-normal-mle}
设 $X_i\sim N(\mu,\sigma^2)$，$\mu,\sigma^2$ 均未知，求最大似然估计。
\end{example}

\begin{solve}
对给定样本，令 $\overline x=n^{-1}\sum x_i$、$Q=\sum(x_i-\overline x)^2$。对数似然为
\[
  \ell(\mu,\sigma^2)
   =-\frac n2\log(2\pi)-\frac n2\log\sigma^2
      -\frac1{2\sigma^2}\sum_{i=1}^n(x_i-\mu)^2.
\]
利用 $\sum(x_i-\mu)^2=Q+n(\overline x-\mu)^2$，固定 $\sigma^2$ 时唯一的最优位置为 $\mu=\overline x$。代入后令 $v=\sigma^2>0$，有
\[
  \frac{d}{dv}\ell(\overline x,v)=\frac{Q-nv}{2v^2}.
\]
若 $Q>0$，导数在 $v=Q/n$ 左侧为正、右侧为负，故得到全局最大值
\begin{equation}
  \widehat\mu_{\mathrm{ML}}=\overline X,\qquad
  \widehat{\sigma^2}_{\mathrm{ML}}=M_2.
  \label{eq:study-normal-mle}
\end{equation}
当 $n\ge2$ 时，$Q>0$ 在正态抽样下几乎必然成立。若所有观测完全相同，$Q=0$，则令 $\mu=\overline x,\sigma^2\downarrow0$ 会使似然无界，在 $\sigma^2>0$ 内没有最大值；$n=1$ 必属此情形。若均值 $\mu_0$ 已知，则方差的 MLE 改为 $n^{-1}\sum(X_i-\mu_0)^2$，不能再减去 $\overline X$。
\end{solve}

对于指数速率，$\ell(\lambda)=n\log\lambda-\lambda\sum x_i$，在 $\sum x_i>0$ 时唯一最大值点为 $\widehat\lambda_{\mathrm{ML}}=1/\overline X$。对于泊松参数，$\ell(\lambda)=-n\lambda+(\sum x_i)\log\lambda-\sum\log(x_i!)$，在 $\lambda\ge0$ 的参数空间内得到 $\widehat\lambda_{\mathrm{ML}}=\overline X$，全零样本对应 $\lambda=0$。若模型要求 $\lambda>0$，全零样本只有在 $\lambda\downarrow0$ 时达到上确界，MLE 不存在。二项总体 $\operatorname{Bin}(N,p)$ 在 $p\in[0,1]$ 内的 MLE 为 $\overline X/N$，全失败和全成功样本分别在 $0,1$ 达到最大；若把参数空间限定为 $(0,1)$，同样要单独处理这两种情形。

\begin{example}
\label{ex:study-uniform-mle}
设 $X_i\sim U(0,\theta)$，$\theta>0$，求最大似然估计。
\end{example}

\begin{solve}
采用端点包含在支撑中的密度版本 $f(x;\theta)=\theta^{-1}\mathbf1_{\{0\le x\le\theta\}}$。对正的观测值，有
\begin{equation}
  L(\theta;\boldsymbol x)=\theta^{-n}
      \mathbf1_{\{\theta\ge x_{(n)}\}},
  \qquad
  \widehat\theta_{\mathrm{ML}}=X_{(n)}.
  \label{eq:study-uniform-mle}
\end{equation}
当 $\theta<x_{(n)}$ 时，参数不能解释当前样本；在 $\theta\ge x_{(n)}$ 的可行域内，$\theta^{-n}$ 严格递减，因此最大值位于边界。若直接删除示性函数再求导，会丢失决定答案的约束。连续分布的端点取值不改变其概率分布，但会影响逐点似然是否在边界取到最大；这里明确端点约定，以得到教材采用的 MLE。
\end{solve}

\begin{example}
\label{ex:study-uniform-location-mle}
设 $X_i\sim U(\theta-\tfrac12,\theta+\tfrac12)$，求位置参数的最大似然估计。
\end{example}

\begin{solve}
仍取包含端点的密度版本，则
\[
  L(\theta;\boldsymbol x)
    =\mathbf1_{\{x_{(n)}-1/2\le\theta\le x_{(1)}+1/2\}}.
\]
因此区间 $[x_{(n)}-1/2,x_{(1)}+1/2]$ 中每一点都是最大值点；若另有 $\theta>0$ 的约束，则与 $(0,\infty)$ 取交集。在来自该模型的有限样本下，交集几乎必然非空。中程数 $(X_{(1)}+X_{(n)})/2$ 是无额外约束时的一种选取。这说明最大似然估计不一定唯一，也不一定来自似然方程。
\end{solve}

类似地，当 $U(a,b)$ 的两个端点均未知且样本极差为正时，似然为 $(b-a)^{-n}\mathbf1_{\{a\le x_{(1)},\,b\ge x_{(n)}\}}$，最小的可行区间给出 $\widehat a_{\mathrm{ML}}=X_{(1)}$、$\widehat b_{\mathrm{ML}}=X_{(n)}$。若极差为零，则可以令区间长度趋于零，似然无界，不能使用这一结论。

\subsection{离散参数与不存在矩的总体}

\begin{example}
\label{ex:study-capture-mle}
教材例 6.7 中，先标记并放回 $500$ 尾鱼，再随机捕获 $400$ 尾，其中标记鱼有 $60$ 尾。假定总体封闭、标记不脱落，标记鱼充分混合且与其他鱼具有相同捕获机会，求总数 $N$ 的最大似然估计。
\end{example}

\begin{solve}
第二次捕获数中的标记鱼数服从超几何分布。固定观测结果后，似然为
\[
  L(N)=\frac{\binom{500}{60}\binom{N-500}{340}}{\binom N{400}},
  \qquad N=840,841,\ldots.
\]
对 $N\ge841$，相邻似然比为
\[
  \frac{L(N)}{L(N-1)}
    =\frac{(N-500)(N-400)}{N(N-840)}.
\]
分子减去分母等于 $200000-60N$，所以似然先增后减，在 $N=\lfloor200000/60\rfloor=3333$ 达到最大。这里的 $N$ 是整数参数，不需要对它求导。一般地，若对应阈值恰为可行域内部的整数 $r$，相邻比在 $N=r$ 等于 $1$，则 $r-1,r$ 都是最大值点，不能忽略并列最大值。
\end{solve}

\begin{example}
\label{ex:study-cauchy-location}
设总体具有柯西位置分布（Cauchy location distribution）
\[
  f(x;\theta)=\frac1{\pi[1+(x-\theta)^2]},
  \qquad x,\theta\in\mathbb R.
\]
讨论位置参数的估计。
\end{example}

\begin{solve}
柯西分布没有有限的一阶绝对矩，因而不能用 $\mathbb E(X)=\theta$ 建立通常的一阶矩方程。最大似然法仍可使用，其对数似然和似然方程为
\[
  \ell(\theta)=-n\log\pi-\sum_{i=1}^n\log[1+(x_i-\theta)^2],
  \qquad
  \ell'(\theta)=2\sum_{i=1}^n\frac{x_i-\theta}{1+(x_i-\theta)^2}=0.
\]
由于 $\ell$ 连续且在 $|\theta|\to\infty$ 时趋于 $-\infty$，最大值存在，但似然方程可能有多个根，求得一个根后仍须比较似然值。另一方面，$\theta$ 是该分布唯一的中位数，可直接用样本中位数
\[
  \widetilde X=
  \begin{cases}
    X_{((n+1)/2)},&n\text{ 为奇数},\\
    (X_{(n/2)}+X_{(n/2+1)})/2,&n\text{ 为偶数}
  \end{cases}
\]
来估计。它是利用分位数构造的估计，不应仅因形式简单就称为最大似然估计；其优劣仍需在给定准则下研究。
\end{solve}

\section{优良性准则}
\label{sec:study-estimator-criteria}

评价估计量时应把样本重新视为随机向量。某一次估计碰巧接近真值，并不能说明该方法在重复抽样中更准确；一种方法在某个参数值下较好，也不意味着它在整个参数空间内都较好。以下先考虑标量目标 $g(\theta)$，所有涉及期望或方差的结论均要求相应量存在。

\subsection{偏差、无偏性及修正}

\begin{definition}[偏差与无偏性]
\label{def:study-estimator-bias}
估计量 $T$ 的\emph{偏差（bias）}为
\begin{equation}
  b_\theta(T)=\mathbb E_\theta(T)-g(\theta).
  \label{eq:study-estimator-bias}
\end{equation}
若对每个 $\theta\in\Theta$ 都有 $b_\theta(T)=0$，则称 $T$ 为 $g(\theta)$ 的\emph{无偏估计量（unbiased estimator）}。若估计量序列 $T_n$ 满足对每个 $\theta$ 都有 $\mathbb E_\theta(T_n)\to g(\theta)$，则称其为渐近无偏估计。
\end{definition}

无偏性是抽样分布的均值性质，不是每次观测都估计正确。固定样本量 $n$，独立重复整个抽样过程，若 $\mathbb E_\theta|T|<\infty$，则这些重复估计值的平均依强大数定律趋于 $\mathbb E_\theta T$；无偏时这个极限恰好是 $g(\theta)$。这里增加的是独立重复抽样的次数，不能与单次样本量 $n\to\infty$ 混为一谈。

\begin{proposition}[均值和方差的无偏估计]
\label{prop:study-unbiased-mean-variance}
若 $X_i$ 独立同分布，$\mathbb E(X_i)=\mu$、$\operatorname{Var}(X_i)=\sigma^2<\infty$，则
\begin{equation}
  \mathbb E(\overline X)=\mu,\qquad
  \operatorname{Var}(\overline X)=\frac{\sigma^2}{n},\qquad
  \mathbb E(M_2)=\frac{n-1}{n}\sigma^2.
  \label{eq:study-mean-variance-expectations}
\end{equation}
因此 $n\ge2$ 时，式~\eqref{eq:study-unbiased-sample-variance} 中的 $S^2$ 是总体方差的无偏估计；结论不要求总体正态。
\end{proposition}

\begin{proof}
期望的线性性与样本独立性给出前两式。再由恒等式
\[
  \sum_{i=1}^n(X_i-\overline X)^2
    =\sum_{i=1}^n(X_i-\mu)^2-n(\overline X-\mu)^2,
\]
取期望得到 $n\mathbb E(M_2)=n\sigma^2-n(\sigma^2/n)=(n-1)\sigma^2$。乘以 $1/(n-1)$ 即得 $\mathbb E(S^2)=\sigma^2$。
\end{proof}

残差满足 $\sum(X_i-\overline X)=0$，因而落在一个 $n-1$ 维线性子空间中，这给出自由度减少一的几何解释，但不能据此声称一般总体的各个残差相互独立。无偏性也不在非线性变换下保留：由柯西--施瓦茨不等式，$\mathbb E(S)\le\sqrt{\mathbb E(S^2)}=\sigma$，当 $S$ 非退化时严格小于 $\sigma$。所以 $S^2$ 无偏并不意味着 $S$ 无偏；同理，$\overline X$ 对 $1/\lambda$ 无偏也不意味着 $1/\overline X$ 对 $\lambda$ 无偏。

\begin{example}
\label{ex:study-uniform-unbiased-correction}
设 $X_i\sim U(0,\theta)$，计算最大值的偏差并作无偏修正。
\end{example}

\begin{solve}
对 $0\le x\le\theta$，独立性给出
\[
  \mathbb P_\theta(X_{(n)}\le x)=\left(\frac x\theta\right)^n,
  \qquad
  f_{X_{(n)}}(x)=\frac{nx^{n-1}}{\theta^n},\quad0<x<\theta.
\]
因此，对 $r>0$，
\begin{equation}
  \mathbb E_\theta(X_{(n)}^r)
    =\int_0^\theta x^r\frac{nx^{n-1}}{\theta^n}\,dx
    =\frac n{n+r}\theta^r.
  \label{eq:study-uniform-maximum-moments}
\end{equation}
特别地，$\mathbb E_\theta(X_{(n)})=n\theta/(n+1)$，偏差为 $-\theta/(n+1)$，故
\begin{equation}
  \widetilde\theta=\frac{n+1}{n}X_{(n)},\qquad
  \operatorname{Var}_\theta(\widetilde\theta)
     =\frac{\theta^2}{n(n+2)}
  \label{eq:study-uniform-corrected-maximum}
\end{equation}
是无偏修正及其方差。修正后的统计量不再是原来的 MLE；原来的 MLE 虽有偏，却是渐近无偏的。
\end{solve}

\begin{example}
\label{ex:study-no-unbiased-reciprocal}
设只观测 $Y\sim\operatorname{Bin}(N,p)$，其中正整数 $N$ 已知、$0<p<1$。证明 $1/p$ 不存在实值无偏估计量。
\end{example}

\begin{solve}
任一实值统计量由有限个数 $t_0,\ldots,t_N$ 决定，其期望
\[
  P(p)=\sum_{j=0}^Nt_j\binom Njp^j(1-p)^{N-j}
\]
是次数至多为 $N$ 的多项式。若它恒等于 $1/p$，则多项式 $pP(p)-1$ 在 $(0,1)$ 的每一点为零，故必须为零多项式；但其在 $p=0$ 处的值为 $-1$，矛盾。无偏估计不一定存在，因此无偏性不能成为不加区别的唯一要求。
\end{solve}

\subsection{均方误差、有效性与最小方差无偏估计}

\begin{definition}[均方误差]
\label{def:study-estimator-mse}
若 $\mathbb E_\theta(T^2)<\infty$，定义估计 $g(\theta)$ 的\emph{均方误差（mean squared error，MSE）}为
\begin{equation}
  \operatorname{MSE}_\theta(T)
     =\mathbb E_\theta[(T-g(\theta))^2]
     =\operatorname{Var}_\theta(T)+b_\theta(T)^2.
  \label{eq:study-mse-decomposition}
\end{equation}
\end{definition}

最后一个等式由 $T-g=(T-\mathbb E_\theta T)+b_\theta(T)$ 展开平方并取期望得到，其中交叉项为零。它把随机波动与系统偏差纳入同一个准则。对无偏估计，MSE 等于方差；对有偏估计，只比较方差可能得出错误结论。也可以采用平均绝对误差 $\mathbb E_\theta|T-g(\theta)|$ 等准则，不同损失函数对应不同的评价目标。

\begin{definition}[有效性与最小方差无偏估计]
\label{def:study-mvue}
若 $T_1,T_2$ 都是 $g(\theta)$ 的无偏估计，且对所有 $\theta\in\Theta$ 都有
\[
  \operatorname{Var}_\theta(T_1)\le\operatorname{Var}_\theta(T_2),
\]
并至少在一个参数值处严格小于，则称 $T_1$ 比 $T_2$ 更有效。若无偏估计 $T_*$ 对每个具有有限方差的无偏估计 $T$ 都满足
\begin{equation}
  \operatorname{Var}_\theta(T_*)\le\operatorname{Var}_\theta(T),
  \qquad\theta\in\Theta,
  \label{eq:study-umvue}
\end{equation}
则称 $T_*$ 为\emph{一致最小方差无偏估计量（uniformly minimum variance unbiased estimator，UMVUE）}，教材简称 MVUE。“一致”指同一个估计量在整个参数空间逐点最优，不是只在某个未知真值处最优。
\end{definition}

将一个估计量与少数候选量作比较，只能得到这些候选量之间的结论，不能据此宣布它在所有无偏估计中最优。例如，设总体均值为 $\mu$、方差为 $\sigma^2>0$，取不含未知参数的常数权重 $\sum_{i=1}^nw_i=1$，则 $T_w=\sum w_iX_i$ 无偏，且
\[
  \operatorname{Var}(T_w)=\sigma^2\sum_{i=1}^nw_i^2
    \ge\frac{\sigma^2}{n}
    =\operatorname{Var}(\overline X).
\]
这里使用 $(\sum w_i)^2\le n\sum w_i^2$，等号当且仅当 $w_i=1/n$。这证明样本均值在常系数线性无偏估计类内最优，尚未覆盖所有非线性无偏估计。

在均匀总体 $U(0,\theta)$ 中，矩估计 $2\overline X$ 和式~\eqref{eq:study-uniform-corrected-maximum} 中的 $\widetilde\theta$ 都无偏，而
\[
  \operatorname{Var}_\theta(2\overline X)=\frac{\theta^2}{3n},
  \qquad
  \operatorname{Var}_\theta(\widetilde\theta)
       =\frac{\theta^2}{n(n+2)}.
\]
当 $n>1$ 时后者严格较小，说明即使样本均值在线性类内最优，使用次序统计量仍可能进一步降低方差。无偏修正也不必降低 MSE：方差与偏差平方的权衡见本章习题中教材第 14 题的解答。

\subsection{克拉默--拉奥方差下界}

要证明某个估计量在所有无偏估计中最优，可以先寻找共同的方差下界。下面只讨论一维参数 $\theta$，并明确正则条件：参数空间为开区间，样本支撑与参数无关，密度或概率质量函数关于参数可微，且对归一化条件和下述无偏期望求导时均可交换微分与积分（离散情形为求和）。例如，在每个参数邻域内，相关导数被可积函数控制，是保证这些交换的一组充分条件。此外要求信息量正且有限、估计量方差有限。

\begin{definition}[得分与费希尔信息]
\label{def:study-fisher-information}
单个观测的\emph{得分（score）}和\emph{费希尔信息（Fisher information）}定义为
\begin{equation}
  u_\theta(X)=\frac{\partial}{\partial\theta}\log f(X;\theta),
  \qquad
  I_1(\theta)=\mathbb E_\theta[u_\theta(X)^2].
  \label{eq:study-fisher-information}
\end{equation}
样本得分为 $U_n(\theta)=\sum_{i=1}^nu_\theta(X_i)$。在上述正则条件下，$\mathbb E_\theta u_\theta=0$，由独立性有
\[
  I_n(\theta):=\operatorname{Var}_\theta(U_n(\theta))
      =nI_1(\theta).
\]
\end{definition}

\begin{theorem}[克拉默--拉奥不等式]
\label{thm:study-cramer-rao}
在上述正则条件下，设 $T$ 是可微参数函数 $g(\theta)$ 的无偏估计，且允许在 $\mathbb E_\theta T=g(\theta)$ 中交换微分与积分，则
\begin{equation}
  \operatorname{Var}_\theta(T)
    \ge\frac{[g'(\theta)]^2}{nI_1(\theta)}.
  \label{eq:study-cramer-rao}
\end{equation}
在给定 $\theta$ 处等号成立，当且仅当
\begin{equation}
  T-g(\theta)
     =\frac{g'(\theta)}{nI_1(\theta)}U_n(\theta)
     \quad\mathbb P_\theta\text{-几乎处处}.
  \label{eq:study-cramer-rao-equality}
\end{equation}
\end{theorem}

\begin{proof}
设联合密度为 $f_n(\boldsymbol x;\theta)$。对 $\int f_n\,d\boldsymbol x=1$ 和 $\int T f_n\,d\boldsymbol x=g(\theta)$ 求导，利用 $\partial_\theta f_n=U_n f_n$，得到
\[
  \mathbb E_\theta U_n=0,\qquad
  \mathbb E_\theta(TU_n)=g'(\theta).
\]
故 $\operatorname{Cov}_\theta(T,U_n)=g'(\theta)$。协方差的柯西--施瓦茨不等式给出
\[
  [g'(\theta)]^2
    \le\operatorname{Var}_\theta(T)
       \operatorname{Var}_\theta(U_n)
    =\operatorname{Var}_\theta(T)\,nI_1(\theta),
\]
除以正数 $nI_1(\theta)$ 即得下界。等号条件也可由
\[
  \mathbb E_\theta\!\left[
    \left(T-g(\theta)-\frac{g'(\theta)}{nI_1(\theta)}U_n(\theta)\right)^2
  \right]
   =\operatorname{Var}_\theta(T)
     -\frac{[g'(\theta)]^2}{nI_1(\theta)}
\]
看出：右端为零，当且仅当括号内的随机变量几乎处处为零。
\end{proof}

从几何上看，应在零均值、平方可积的随机变量空间中，以 $\mathbb E_\theta(ZY)$ 为内积，将 $T-g(\theta)$ 投影到得分方向；方差下界就是投影长度的平方。这样可避免把所有随机变量上的协方差直接当作内积，因为非零常数的协方差范数为零。下界不一定能够达到；若某个无偏估计对所有参数都达到下界，它必为 UMVUE，反过来 UMVUE 不一定达到该下界。

\begin{example}
\label{ex:study-cramer-rao-applications}
用克拉默--拉奥下界判断正态均值、指数均值和二项比例的样本均值型估计是否为 UMVUE。
\end{example}

\begin{solve}
正态总体 $N(\mu,\sigma^2)$ 中先设 $\sigma^2>0$ 已知，以 $\mu$ 为参数，则
\[
  u_\mu(X)=\frac{X-\mu}{\sigma^2},\qquad
  I_1(\mu)=\frac1{\sigma^2}.
\]
$\mu$ 的无偏估计方差不小于 $\sigma^2/n$，而 $\overline X$ 恰好达到，故为 UMVUE。即使 $\sigma^2$ 也未知，任何对所有 $(\mu,\sigma^2)$ 无偏的估计，在固定任意一个 $\sigma^2$ 的子模型中仍须满足该下界；$\overline X$ 因而仍为 $\mu$ 的 UMVUE。若反而是 $\mu=\mu_0$ 已知，以 $v=\sigma^2$ 为参数，则 $I_1(v)=1/(2v^2)$，无偏估计 $n^{-1}\sum(X_i-\mu_0)^2$ 的方差为 $2v^2/n$，也达到下界。

指数分布 $\operatorname{Exp}(\lambda)$ 中，$u_\lambda(X)=1/\lambda-X$、$I_1(\lambda)=1/\lambda^2$。估计的是均值 $g(\lambda)=1/\lambda$，其导数为 $-1/\lambda^2$，故下界为
\[
  \frac{[g'(\lambda)]^2}{nI_1(\lambda)}
    =\frac1{n\lambda^2}
    =\operatorname{Var}_\lambda(\overline X).
\]
因此 $\overline X$ 是 $1/\lambda$ 的 UMVUE。这一结论不能直接改写成 $1/\overline X$ 是 $\lambda$ 的无偏或最优估计。

对于 $\operatorname{Bin}(N,p)$，$N$ 已知、$0<p<1$，有
\[
  u_p(X)=\frac{X-Np}{p(1-p)},\qquad
  I_1(p)=\frac{N}{p(1-p)}.
\]
估计 $p$ 的方差下界为 $p(1-p)/(nN)$，正好等于无偏估计 $\overline X/N$ 的方差，故后者为 UMVUE。类似地，泊松总体有 $I_1(\lambda)=1/\lambda$，$\overline X$ 以方差 $\lambda/n$ 达到下界；这也证明它优于其他无偏估计，包括用 $S^2$ 估计同一个参数 $\lambda$ 的做法。
\end{solve}

\begin{notealert}[type=caution,title={支撑依赖参数时不能机械套用下界}]
对于 $U(0,\theta)$，若只在支撑内部计算 $\partial_\theta\log f=-1/\theta$，会得到一个期望不为零的“得分”，说明正则推导已失效。积分上限依赖 $\theta$，归一化条件求导时存在边界项。因此不能把 $I_1(\theta)=1/\theta^2$ 代入式~\eqref{eq:study-cramer-rao}，再断言方差至少为 $\theta^2/n$；式~\eqref{eq:study-uniform-corrected-maximum} 的方差实际更小，并不与定理矛盾。
\end{notealert}

\begin{proposition}[均匀总体中修正最大值的最优性]
\label{prop:study-uniform-maximum-umvue}
在 $U(0,\theta)$ 的简单随机样本中，$\widetilde\theta=(n+1)X_{(n)}/n$ 是 $\theta$ 的 UMVUE。
\end{proposition}

\begin{proof}
当 $n=1$ 时也可使用以下论证。记 $M=X_{(n)}$。在给定 $M=m$ 后，最大值所在位置等可能，其余 $n-1$ 个观测的条件分布为独立的 $U(0,m)$，不含 $\theta$。因此，对任意无偏估计 $T$，条件期望 $H(M)=\mathbb E_\theta(T\mid M)$ 可取为不含未知参数的统计量，并且 $\mathbb E_\theta H(M)=\theta$。由全方差公式，$\operatorname{Var}_\theta H(M)\le\operatorname{Var}_\theta T$。

只需证明作为 $M$ 的函数的无偏估计唯一。若 $h(M)$ 在所有 $\theta>0$ 下期望为零，则
\[
  \int_0^\theta h(m)m^{n-1}\,dm=0,\qquad\theta>0.
\]
因 $h(M)$ 可积，被积函数局部可积，对上式按积分上限求导，得到 $h(m)m^{n-1}=0$ 对几乎所有 $m>0$ 成立。于是 $h(M)=0$ 在每个 $\theta$ 下几乎必然成立。取 $h(M)=H(M)-(n+1)M/n$，便知 $H(M)=\widetilde\theta$，再结合方差不等式得到结论。
\end{proof}

\section{点估计量的大样本理论}
\label{sec:study-estimator-asymptotics}

有限样本性质是在固定 $n$ 下研究估计量的抽样分布；大样本性质则研究估计量序列在 $n\to\infty$ 时的行为。两者的区别是研究方式，而不是某个统一的样本量门槛。无偏性控制期望，相合性控制误差概率，渐近正态性进一步描述适当缩放后的误差分布，三者不能互相替代。

\subsection{相合性}

\begin{definition}[相合估计]
\label{def:study-consistent-estimator}
若对每个 $\theta\in\Theta$、每个 $\varepsilon>0$，都有
\begin{equation}
  \mathbb P_\theta(|T_n-g(\theta)|>\varepsilon)\longrightarrow0,
  \qquad n\to\infty,
  \label{eq:study-estimator-consistency}
\end{equation}
则称 $T_n$ 为 $g(\theta)$ 的\emph{弱相合估计量（weakly consistent estimator）}，记为 $T_n\xrightarrow{\mathbb P_\theta}g(\theta)$，简称相合估计。
\end{definition}

定义中的“对每个 $\theta$”是逐点要求，不等于更强的 $\sup_{\theta\in\Theta}\mathbb P_\theta(|T_n-g(\theta)|>\varepsilon)\to0$。依概率收敛也不表示每一条样本路径的误差都单调减小；它表示固定精度以外的误差越来越少见。

\begin{proposition}[由均方误差趋于零判断相合性]
\label{prop:study-mse-consistency}
若对每个 $\theta$，$\operatorname{MSE}_\theta(T_n)\to0$，则 $T_n$ 相合。特别地，渐近无偏且方差趋于零是相合性的一组充分条件。
\end{proposition}

\begin{proof}
由非负随机变量的马尔可夫不等式，
\[
  \mathbb P_\theta(|T_n-g(\theta)|>\varepsilon)
    \le\frac{\mathbb E_\theta[(T_n-g(\theta))^2]}{\varepsilon^2}
    \longrightarrow0.
\]
后一结论再用式~\eqref{eq:study-mse-decomposition} 即得。
\end{proof}

相合性不要求每个有限样本量下都无偏。反过来，无偏也不保证相合：若总体非退化，始终使用 $T_n=X_1$ 估计均值，则虽对每个 $n$ 都无偏，却存在 $\varepsilon>0$ 使 $\mathbb P(|X_1-\mu|>\varepsilon)>0$ 不随 $n$ 改变。

\begin{example}
\label{ex:study-moment-consistency}
设 $X_i$ 独立同分布且 $\mathbb E(X_i^2)<\infty$，证明 $\overline X$、$M_2$ 和 $S^2$ 分别相合地估计均值与方差。
\end{example}

\begin{solve}
由大数定律，
\[
  \overline X\xrightarrow{\mathbb P}\mu,\qquad
  A_2\xrightarrow{\mathbb P}\mathbb E(X^2).
\]
连续函数保持依概率收敛，故
\[
  M_2=A_2-\overline X^2
      \xrightarrow{\mathbb P}\mathbb E(X^2)-\mu^2=\sigma^2.
\]
再由 $S^2=nM_2/(n-1)$ 得 $S^2\xrightarrow{\mathbb P}\sigma^2$。这里不需要四阶矩存在：证明使用二阶原点矩的大数定律，而不是先计算样本方差自身的方差。更一般地，若矩方程在真值附近有连续的逆，且所选样本解落在对应分支上，则由样本矩的收敛可推出该矩估计相合；不能略去解的存在、分支选择等条件。
\end{solve}

\begin{example}
\label{ex:study-uniform-maximum-consistency}
设 $X_i\sim U(0,\theta)$，证明 MLE $X_{(n)}$ 相合，并考察其误差缩放。
\end{example}

\begin{solve}
对 $0<\varepsilon<\theta$，因 $X_{(n)}\le\theta$ 几乎必然成立，
\begin{equation}
  \mathbb P_\theta(|X_{(n)}-\theta|>\varepsilon)
    =\left(1-\frac{\varepsilon}{\theta}\right)^n\longrightarrow0.
  \label{eq:study-uniform-maximum-consistency}
\end{equation}
当 $\varepsilon\ge\theta$ 时该概率为零，故相合性成立。进一步，对固定 $t\ge0$、充分大的 $n$，
\[
  \mathbb P_\theta\!\left(
      \frac{n(\theta-X_{(n)})}{\theta}>t
  \right)=\left(1-\frac tn\right)^n\longrightarrow e^{-t}.
\]
因而 $n(\theta-X_{(n)})/\theta$ 依分布趋于速率为 $1$ 的指数分布。它的自然误差尺度为 $1/n$，极限分布不是正态分布；这也是不能对支撑依赖参数的模型套用一般正则 MLE 结论的例子。
\end{solve}

\subsection{渐近正态性与误差近似}

\begin{definition}[标准化误差的渐近正态性]
\label{def:study-asymptotic-normality}
沿用教材的标准化形式，设 $v_n(\theta)=\operatorname{Var}_\theta(T_n)\in(0,\infty)$。若
\begin{equation}
  \frac{T_n-g(\theta)}{\sqrt{v_n(\theta)}}
      \xrightarrow{\mathcal D}N(0,1),
  \label{eq:study-asymptotic-normality}
\end{equation}
则称该标准化误差具有\emph{渐近正态性（asymptotic normality）}，其中 $\xrightarrow{\mathcal D}$ 表示依分布收敛。
\end{definition}

常用的另一种表达是 $\sqrt n(T_n-g(\theta))\xrightarrow{\mathcal D}N(0,v(\theta))$，其中 $0<v(\theta)<\infty$。若另有 $nv_n(\theta)\to v(\theta)$，它与式~\eqref{eq:study-asymptotic-normality} 等价；若尚未证明二阶矩收敛，就不能仅凭分布极限把 $v(\theta)$ 当作 $\lim n\operatorname{Var}_\theta(T_n)$。正态近似通常用于相合估计，在上述标准化形式下还应有 $v_n(\theta)\to0$ 才能据此推出误差趋于零。

\begin{theorem}[样本均值及其可行标准化]
\label{thm:study-studentized-mean-limit}
若 $X_i$ 独立同分布，均值为 $\mu$，且 $0<\sigma^2<\infty$，则
\begin{equation}
  \frac{\sqrt n(\overline X-\mu)}{\sigma}
       \xrightarrow{\mathcal D}N(0,1),\qquad
  \frac{\sqrt n(\overline X-\mu)}{S}
       \xrightarrow{\mathcal D}N(0,1).
  \label{eq:study-studentized-mean-limit}
\end{equation}
第二个统计量在 $S=0$ 时可任意约定其值，该事件的概率随 $n$ 增大趋于零。
\end{theorem}

\begin{proof}
第一式是独立同分布中心极限定理。由例~\ref{ex:study-moment-consistency}，$S^2\xrightarrow{\mathbb P}\sigma^2>0$，故 $S/\sigma\xrightarrow{\mathbb P}1$，且 $\mathbb P(S=0)\to0$。将第一式的左端乘以 $\sigma/S$，由斯卢茨基定理（Slutsky's theorem）得到第二式：依分布收敛的随机变量与依概率收敛到常数的随机变量相乘，其极限为相应极限之积。
\end{proof}

因此，在大样本下 $\overline X$ 的标准误差可估为 $S/\sqrt n$。总体本身不必正态；只有总体正态时，固定有限 $n$ 的 $(\overline X-\mu)/(S/\sqrt n)$ 才具有精确的 $t_{n-1}$ 分布。大样本近似的精度取决于总体分布，不能只凭一个固定样本量阈值保证。

参数函数的极限可用\emph{德尔塔方法（delta method）}得到。若 $\widehat\theta_n\xrightarrow{\mathbb P}\theta$、$\sqrt n(\widehat\theta_n-\theta)\xrightarrow{\mathcal D}N(0,v)$，且 $g$ 在 $\theta$ 处可微，则一阶展开给出
\begin{equation}
  \sqrt n\,[g(\widehat\theta_n)-g(\theta)]
      \xrightarrow{\mathcal D}
          N\!\left(0,[g'(\theta)]^2v\right).
  \label{eq:study-delta-method}
\end{equation}
其依据是可微性将差值写成 $[g'(\theta)+r_n](\widehat\theta_n-\theta)$，其中 $r_n\xrightarrow{\mathbb P}0$，再用斯卢茨基定理。若 $g'(\theta)=0$，右端退化为零，不能继续按非退化正态分布标准化，往往需要更高阶展开。

例如，指数分布下 $\mathbb E X=1/\lambda$、$\operatorname{Var}X=1/\lambda^2$，对 $g(t)=1/t$ 使用式~\eqref{eq:study-delta-method}，得到
\[
  \sqrt n\left(\frac1{\overline X}-\lambda\right)
      \xrightarrow{\mathcal D}N(0,\lambda^2).
\]
泊松总体下，零次事件概率 $q=e^{-\lambda}$ 的 MLE 为 $\widehat q=e^{-\overline X}$；当 $\lambda>0$ 时，同理有
\[
  \sqrt n(\widehat q-q)
      \xrightarrow{\mathcal D}N(0,\lambda e^{-2\lambda}).
\]
这些是对具体模型的直接验证。一般 MLE 的相合性与渐近正态性还涉及参数可识别、真值为内点、似然的光滑性与信息量等条件，不是最大化定义本身的无条件推论。

\section{扩展阅读：德军坦克问题}
\label{sec:study-german-tank}

将教材中的序列号问题抽象为：总体编号为 $1,2,\ldots,N$，其中 $N$ 为未知正整数；从中等可能、不放回地抽取 $k$ 个不同编号，$k$ 已知且 $1\le k\le N$。记抽到的编号集合为 $A$，最大编号为 $M=\max A$，其观测值为 $m$。关键假设是每个大小为 $k$ 的编号子集等可能；若某些编号更容易被观测，以下估计的性质便可能改变。这是有限总体的不放回抽样，不能把编号视为独立的连续均匀样本。

\subsection{最大编号的分布及无偏修正}

在 $\binom Nk$ 个等可能子集中，最大值等于 $m$ 时必须包含 $m$，其余 $k-1$ 个编号从 $1,\ldots,m-1$ 中选择，故
\begin{equation}
  \mathbb P_N(M=m)
    =\frac{\binom{m-1}{k-1}}{\binom Nk},
  \qquad m=k,k+1,\ldots,N.
  \label{eq:study-tank-maximum-mass}
\end{equation}
利用组合恒等式 $m\binom{m-1}{k-1}=k\binom mk$ 及 $\sum_{m=k}^N\binom mk=\binom{N+1}{k+1}$，有
\begin{equation}
  \mathbb E_N M
    =\frac{k\binom{N+1}{k+1}}{\binom Nk}
    =\frac{k(N+1)}{k+1}.
  \label{eq:study-tank-maximum-mean}
\end{equation}
因此，直接用 $M$ 会低估 $N$，偏差为 $-(N-k)/(k+1)$；解出期望关系并用 $M$ 替换 $\mathbb E_NM$，得到
\begin{equation}
  \widehat N_{\mathrm U}=\frac{k+1}{k}M-1,
  \qquad \mathbb E_N\widehat N_{\mathrm U}=N.
  \label{eq:study-tank-unbiased-estimator}
\end{equation}
它可解释为在最大观测编号之后补上一个估计的平均间隔。该估计量可能不是整数；若再四舍五入，应另行评价舍入后的偏差，不能沿用原来的无偏结论。

为了评价精度，进一步利用
\[
  m(m+1)\binom{m-1}{k-1}
      =k(k+1)\binom{m+1}{k+1}
\]
求和，得到
\[
  \mathbb E_N[M(M+1)]=\frac{k(N+1)(N+2)}{k+2}.
\]
减去 $\mathbb E_N M$ 得到二阶矩，再减去均值平方，可得
\begin{equation}
  \operatorname{Var}_N(M)
    =\frac{k(N+1)(N-k)}{(k+1)^2(k+2)},\qquad
  \operatorname{Var}_N(\widehat N_{\mathrm U})
    =\frac{(N+1)(N-k)}{k(k+2)}.
  \label{eq:study-tank-estimator-variance}
\end{equation}
当 $N\gg k$ 时，后者约为 $N^2/[k(k+2)]$，标准差约为 $N/\sqrt{k(k+2)}$；若 $k$ 还较大，可进一步近似为 $N/k$。这里 $N/k$ 是标准差的近似量级，不能与方差混淆。

\subsection{最大似然估计与最小方差性质}

固定观测集合 $A$ 后，
\[
  L(N;A)=\binom Nk^{-1}\mathbf1_{\{N\ge m\}},
  \qquad N=k,k+1,\ldots.
\]
因为 $\binom Nk$ 随 $N\ge k$ 严格递增，故 $\widehat N_{\mathrm{ML}}=M$。最大似然估计与无偏修正不同：前者使当前样本的似然最大，后者使重复抽样的平均值等于总体大小，这两个目标没有理由产生同一个统计量。

\begin{proposition}[坦克总数的最小方差无偏估计]
\label{prop:study-tank-umvue}
在上述等可能不放回抽样模型下，式~\eqref{eq:study-tank-unbiased-estimator} 是 $N$ 的 UMVUE。
\end{proposition}

\begin{proof}
给定 $M=m$ 后，样本必含 $m$，其余元素在 $1,\ldots,m-1$ 的所有 $k-1$ 元子集中等可能，条件分布不依赖 $N$。因此，对任意无偏估计 $T(A)$，$H(M)=\mathbb E_N(T\mid M)$ 是不含 $N$ 的统计量，仍无偏，且由全方差公式有 $\operatorname{Var}_N H(M)\le\operatorname{Var}_N T$。

若一个函数 $h(M)$ 对每个 $N\ge k$ 都满足 $\mathbb E_Nh(M)=0$，则
\[
  \sum_{m=k}^Nh(m)\binom{m-1}{k-1}=0,\qquad N=k,k+1,\ldots.
\]
取 $N=k$ 得 $h(k)=0$；若已知 $h(k),\ldots,h(N-1)$ 都为零，等式便给出 $h(N)\binom{N-1}{k-1}=0$，故 $h(N)=0$。由归纳法，$h$ 在所有可能取值上均为零。因此 $M$ 的函数中无偏估计唯一，必为 $\widehat N_{\mathrm U}$，再结合条件平均的方差不等式即得最优性。
\end{proof}

若从总体一阶矩 $\mathbb E_N X=(N+1)/2$ 出发，也能构造无偏估计 $2\overline X-1$；这与利用最大编号的期望关系作修正是不同的构造。坦克问题说明，次序统计量能有效使用参数决定的支撑信息，但这种优势以抽样模型成立为前提，并不能仅凭某次估计值接近真值来验证模型。

\section*{本章总结}
\addcontentsline{toc}{section}{本章总结}

点估计首先要区分固定未知参数、随机估计量和观测后的估计值。矩估计通过样本矩与总体矩匹配构造参数估计，须检查矩的存在、方程的可解性和参数约束；最大似然估计固定数据后比较参数，须保留支撑条件、检查边界与全局最大值，并允许不唯一或不存在的情形。对参数函数作代入估计很方便，但非线性变换一般不保持无偏性。

无偏性是期望正确，MSE 同时计算方差与偏差平方；有效性是在给定估计类与整个参数空间内的比较。证明 UMVUE 需要与所有无偏估计比较，正则模型中可使用克拉默--拉奥下界，支撑依赖参数的模型则要另作论证。相合性说明误差概率趋于零，渐近正态性说明标准化误差的分布近似；无偏、相合和渐近正态之间均不能无条件互推。均匀总体最大值与坦克编号最大值都能作无偏修正，但前者来自独立连续抽样，后者来自有限总体不放回抽样，其分布和修正公式必须分别推导。

\section*{习题}
\addcontentsline{toc}{section}{习题}

以下选解教材第六章第 12、14、20、27、44、54、56、58 题，分别检验次序统计量的修正、偏差与方差的权衡、加权估计、参数函数估计及相合性；题目中的端点密度沿用本章包含边界的版本，不改变原分布。

\begin{example}
\label{ex:study-ch6-ex12}
教材第 12 题。设 $X_1,X_2,X_3$ 为 $U(0,\theta)$ 的简单随机样本。证明 $\tfrac43X_{(3)}$ 和 $4X_{(1)}$ 都是 $\theta$ 的无偏估计，并比较有效性。
\end{example}

\begin{solve}
由式~\eqref{eq:study-uniform-maximum-moments}，$\mathbb E X_{(3)}=3\theta/4$，$\operatorname{Var}(X_{(3)})=3\theta^2/80$。由于 $\theta-X_i$ 仍为独立的 $U(0,\theta)$ 变量，$X_{(1)}$ 与 $\theta-X_{(3)}$ 同分布，故 $\mathbb E X_{(1)}=\theta/4$，方差同样为 $3\theta^2/80$。于是两个估计均无偏，而
\[
  \operatorname{Var}\!\left(\frac43X_{(3)}\right)
     =\frac{\theta^2}{15},\qquad
  \operatorname{Var}(4X_{(1)})=\frac{3\theta^2}{5}.
\]
前者的方差为后者的 $1/9$，因此更有效；虽然最大值与最小值本身的方差相同，无偏修正所需倍数的不同改变了最终精度。
\end{solve}

\begin{example}
\label{ex:study-ch6-ex14}
教材第 14 题。设 $X_i\sim N(\mu,\sigma^2)$，$n\ge2$，记 $Q=\sum(X_i-\overline X)^2$。比较无偏估计 $Q/(n-1)$ 与有偏估计 $Q/(n+1)$ 对 $\sigma^2$ 的均方误差。
\end{example}

\begin{solve}
对正态样本作正交变换，残差平方和 $Q/\sigma^2$ 可写成 $n-1$ 个独立标准正态变量的平方和，故服从 $\chi^2_{n-1}$，从而
\[
  \mathbb E Q=(n-1)\sigma^2,\qquad
  \operatorname{Var}(Q)=2(n-1)\sigma^4.
\]
在常数倍估计类 $cQ$ 中，
\[
  \frac{\operatorname{MSE}(cQ)}{\sigma^4}
     =2(n-1)c^2+[(n-1)c-1]^2
     =(n-1)(n+1)c^2-2(n-1)c+1.
\]
此二次式在 $c=1/(n+1)$ 处最小，因此
\[
  \operatorname{MSE}\!\left(\frac Q{n+1}\right)
     =\frac{2\sigma^4}{n+1}
     <\frac{2\sigma^4}{n-1}
     =\operatorname{MSE}\!\left(\frac Q{n-1}\right).
\]
$Q/(n+1)$ 的偏差为 $-2\sigma^2/(n+1)$，但方差的下降超过偏差平方的增加。这证明有偏估计可能具有更小的 MSE；这里只在 $cQ$ 这一估计类内求得最优值。
\end{solve}

\begin{example}
\label{ex:study-ch6-ex20}
教材第 20 题。$k$ 台仪器独立测量同一常数 $\theta$，所得 $X_i$ 满足 $\mathbb E X_i=\theta$、$\operatorname{Var}(X_i)=\sigma_i^2$，其中 $\sigma_i^2>0$ 已知。求常系数线性无偏估计 $T=\sum_{i=1}^ka_iX_i$ 中方差最小的权重。
\end{example}

\begin{solve}
在以所有实数 $\theta$ 为可能值的模型中，无偏条件为 $\sum a_i=1$，独立性给出 $\operatorname{Var}(T)=\sum a_i^2\sigma_i^2$。由
\[
  1=\left(\sum_{i=1}^k(a_i\sigma_i)\frac1{\sigma_i}\right)^2
    \le\left(\sum_{i=1}^ka_i^2\sigma_i^2\right)
       \left(\sum_{i=1}^k\sigma_i^{-2}\right),
\]
最小方差及其达到条件为
\[
  a_i=\frac{\sigma_i^{-2}}{\sum_{j=1}^k\sigma_j^{-2}},
  \qquad
  \min\operatorname{Var}(T)
       =\frac1{\sum_{j=1}^k\sigma_j^{-2}}.
\]
即按方差倒数加权，精度越高的仪器权重越大；方差相同时退化为等权平均。结论限于给定的线性无偏类，且依赖测量独立与方差已知。
\end{solve}

\begin{example}
\label{ex:study-ch6-ex27}
教材第 27 题。设 $X_i\sim\operatorname{Pois}(\lambda)$，求 $q=\mathbb P_\lambda(X=0)$ 的 MLE。$122$ 个观测中，取值 $0,1,2,3,4,5$ 的频数依次为 $44,42,21,9,4,2$，没有更大值，求相应估计值。
\end{example}

\begin{solve}
由 $q=e^{-\lambda}$ 及最大似然估计的不变性，
\[
  \widehat q_{\mathrm{ML}}=e^{-\overline X},\qquad
  \overline x=\frac{42+2\cdot21+3\cdot9+4\cdot4+5\cdot2}{122}
     =\frac{137}{122},
\]
故估计值为 $e^{-137/122}\approx0.3253$。它利用泊松模型中的全部频数，不同于直接使用零值频率 $44/122$。而且
\[
  \mathbb E_\lambda(e^{-\overline X})
      =\exp\!\left[n\lambda(e^{-1/n}-1)\right]
\]
一般不等于 $e^{-\lambda}$，所以不变性不意味着变换后仍无偏。
\end{solve}

\begin{example}
\label{ex:study-ch6-ex44}
教材第 44 题。设总体密度为
\[
  f(x;\theta)=\frac1{\sigma}e^{-(x-\theta)/\sigma}
       \mathbf1_{\{x\ge\theta\}},
  \qquad\theta\in\mathbb R,
\]
其中 $\sigma>0$ 已知。求位置参数 $\theta$ 的矩估计、最大似然估计及无偏修正，并比较修正后的有效性。
\end{example}

\begin{solve}
写成 $X=\theta+Y$，其中 $Y$ 是均值为 $\sigma$ 的指数变量，则 $\mathbb E X=\theta+\sigma$、$\operatorname{Var}X=\sigma^2$，故
\[
  \widehat\theta_{\mathrm M}=\overline X-\sigma,\qquad
  \mathbb E\widehat\theta_{\mathrm M}=\theta,\qquad
  \operatorname{Var}(\widehat\theta_{\mathrm M})=\frac{\sigma^2}{n}.
\]
固定样本后的似然为
\[
  L(\theta)=\sigma^{-n}
     \exp\!\left[-\frac{\sum x_i-n\theta}{\sigma}\right]
     \mathbf1_{\{\theta\le x_{(1)}\}}.
\]
它在可行域内随 $\theta$ 严格递增，故 $\widehat\theta_{\mathrm{ML}}=X_{(1)}$。由
\[
  \mathbb P(X_{(1)}-\theta>t)
      =\prod_{i=1}^n\mathbb P(Y_i>t)
      =e^{-nt/\sigma},\qquad t\ge0,
\]
可知 $X_{(1)}-\theta$ 的均值和方差分别为 $\sigma/n,\sigma^2/n^2$。因此无偏修正为 $\widetilde\theta=X_{(1)}-\sigma/n$，方差为 $\sigma^2/n^2$，当 $n>1$ 时比 $\overline X-\sigma$ 更有效。该模型仍属于支撑依赖参数的情形。
\end{solve}

\begin{example}
\label{ex:study-ch6-ex54}
教材第 54 题。设 $X_i\sim\operatorname{Exp}(\lambda)$，采用式~\eqref{eq:study-exponential-rate} 的速率参数化。$\overline X$ 对 $1/\lambda$ 无偏，$1/\overline X$ 是否对 $\lambda$ 无偏？
\end{example}

\begin{solve}
记 $W=\sum_{i=1}^nX_i$，则其密度为
\[
  f_W(w)=\frac{\lambda^n}{\Gamma(n)}
              w^{n-1}e^{-\lambda w}\mathbf1_{\{w>0\}},
\]
即形状为 $n$、速率为 $\lambda$ 的伽马分布。对 $n>1$，直接积分得
\[
  \mathbb E(W^{-1})
    =\frac{\lambda^n}{\Gamma(n)}
       \int_0^\infty w^{n-2}e^{-\lambda w}\,dw
    =\lambda\frac{\Gamma(n-1)}{\Gamma(n)}
    =\frac{\lambda}{n-1}.
\]
因此 $\mathbb E(1/\overline X)=n\lambda/(n-1)>\lambda$，其无偏修正为 $(n-1)/W$。$n=1$ 时倒数的期望发散，不能使用这个修正公式宣称无偏。类似地，当 $n>2$ 时，
\[
  \mathbb E(W^{-2})=\frac{\lambda^2}{(n-1)(n-2)},\qquad
  \operatorname{Var}\!\left(\frac{n-1}{W}\right)
      =\frac{\lambda^2}{n-2}.
\]
在 $n=2$ 时该无偏估计虽有有限期望，方差却无限，进一步说明讨论无偏性与讨论精度所需的矩条件不同。
\end{solve}

\begin{example}
\label{ex:study-ch6-ex56}
教材第 56 题。总体取值 $0,1,2,3$ 的概率依次为 $\theta/2,\theta,3\theta/2,1-3\theta$，取 $\Theta=[0,1/3]$。观测样本为 $(0,3,1,1,0,2,0,0,3,0)$。求矩方程的解与 MLE，讨论无偏性、有效性及参数约束。
\end{example}

\begin{solve}
由 $\mathbb E_\theta X=3-5\theta$，矩方程的形式解为
\[
  T_{\mathrm M}=\frac{3-\overline X}{5}.
\]
记 $N_3=\sum_{i=1}^n\mathbf1_{\{X_i=3\}}$，则忽略与参数无关的正因子后，似然为 $\theta^{n-N_3}(1-3\theta)^{N_3}$。求导并检查两端，可得在闭参数空间上的 MLE
\[
  \widehat\theta_{\mathrm{ML}}=\frac{n-N_3}{3n}.
\]
给定数据中 $n=10$、$\overline x=1$、$N_3=2$，故形式矩解为 $0.4$，MLE 为 $4/15$。\emph[caution]{$0.4>1/3$，所以这组数据的矩方程在参数空间内没有解。}$T_{\mathrm M}$ 仍可作为一个实值点估计量研究，但不能把这个形式解说成满足参数约束的矩估计值。截断到参数空间会改变估计量的性质。

对于未截断的统计量，$\mathbb E_\theta T_{\mathrm M}=\theta$；又因 $n-N_3\sim\operatorname{Bin}(n,3\theta)$，MLE 也无偏。由
\[
  \operatorname{Var}_\theta X
    =(9-20\theta)-(3-5\theta)^2
    =5\theta(2-5\theta)
\]
得到
\[
  \operatorname{Var}_\theta(T_{\mathrm M})
      =\frac{\theta(2-5\theta)}{5n},\qquad
  \operatorname{Var}_\theta(\widehat\theta_{\mathrm{ML}})
      =\frac{\theta(1-3\theta)}{3n}.
\]
两者之差为 $\theta/(15n)$，在 $\theta>0$ 时严格为正，所以 MLE 比未截断的形式矩估计更有效。此例同时说明，无偏性本身不保证估计值始终落在参数空间内。
\end{solve}

\begin{example}
\label{ex:study-ch6-ex58}
教材第 58 题。设 $T_n$ 是 $\theta$ 的无偏估计，$0<\operatorname{Var}_\theta(T_n)<\infty$ 且 $\operatorname{Var}_\theta(T_n)\to0$。证明 $T_n^2$ 是 $\theta^2$ 的相合估计，却不是无偏估计。
\end{example}

\begin{solve}
由无偏性与切比雪夫不等式，
\[
  \mathbb P_\theta(|T_n-\theta|>\varepsilon)
     \le\frac{\operatorname{Var}_\theta(T_n)}{\varepsilon^2}
     \longrightarrow0.
\]
平方函数连续，故 $T_n^2\xrightarrow{\mathbb P}\theta^2$。另一方面，
\[
  \mathbb E_\theta(T_n^2)
      =\theta^2+\operatorname{Var}_\theta(T_n)>\theta^2,
\]
所以每个有限 $n$ 下都不是无偏估计，但其偏差趋于零，因而是渐近无偏的。这里能得到渐近无偏，来自给定的方差条件，并不是依概率收敛本身保证期望收敛。
\end{solve}





